% Impaginazione comune al PDF. Nessun ridimensionamento globale delle tabelle.
\usepackage{needspace}
% Accenti editoriali coordinati con il sito; testo nero su bianco.
\definecolor{libroRosso}{HTML}{9A0000}
\definecolor{libroGrigio}{HTML}{625E59}
\addtokomafont{title}{\color{libroRosso}}
\addtokomafont{part}{\color{libroRosso}}
\addtokomafont{partnumber}{\color{libroGrigio}}
\addtokomafont{chapter}{\color{libroRosso}}
\addtokomafont{chapterprefix}{\color{libroGrigio}}
% Titoli nello stesso carattere del testo (non in sans), come in un libro;
% intestazioni discrete in corsivo grigio, senza filetti.
\setkomafont{disposition}{\normalfont\bfseries}
\addtokomafont{chapter}{\mdseries}
\addtokomafont{part}{\mdseries}
\setkomafont{pageheadfoot}{\normalfont\small\itshape\color{libroGrigio}}
\setkomafont{pagenumber}{\normalfont\small}
% Frontespizio: sottotitolo e affiliazione.
\addtokomafont{subtitle}{\itshape\color{libroGrigio}}
\publishers{\small Dipartimento di Psicologia Generale\\Università di Padova}
% Apertura di capitolo: breve filetto rosso, numero grigio, titolo rosso.
\renewcommand*{\chapterformat}{{\color{libroGrigio}\thechapter}\enskip}
\renewcommand*{\chapterheadstartvskip}{%
  \vspace*{1.5\baselineskip}%
  {\noindent\color{libroRosso}\rule{2.2em}{1.6pt}\par}%
  \vspace{.7\baselineskip}}
% Indice più leggero: capitoli in tondo, parti in grassetto.
\addtokomafont{chapterentry}{\mdseries}
% Riquadri (callout) nei colori del libro invece dell'azzurro predefinito.
\AtBeginDocument{%
  \definecolor{quarto-callout-color}{HTML}{9A0000}%
  \definecolor{quarto-callout-note-color}{HTML}{9A0000}%
  \definecolor{quarto-callout-tip-color}{HTML}{9A0000}%
  \definecolor{quarto-callout-warning-color}{HTML}{9A0000}%
  \definecolor{quarto-callout-important-color}{HTML}{9A0000}%
  \definecolor{quarto-callout-caution-color}{HTML}{9A0000}%
  \definecolor{quarto-callout-color-frame}{HTML}{E3DED8}%
  \definecolor{quarto-callout-note-color-frame}{HTML}{E3DED8}%
  \definecolor{quarto-callout-tip-color-frame}{HTML}{E3DED8}%
  \definecolor{quarto-callout-warning-color-frame}{HTML}{E3DED8}%
  \definecolor{quarto-callout-important-color-frame}{HTML}{E3DED8}%
  \definecolor{quarto-callout-caution-color-frame}{HTML}{E3DED8}%
}
\usepackage{xurl}
\raggedbottom
\setlength{\emergencystretch}{2em}
\widowpenalty=10000
\clubpenalty=10000
\displaywidowpenalty=10000
\renewcommand{\arraystretch}{1.15}
\setlength{\tabcolsep}{5pt}
\RedeclareSectionCommand[beforeskip=0pt,afterskip=16pt]{chapter}
% Le parti restano divisori editoriali intenzionali; niente pagine pari vuote.
\KOMAoptions{open=any}
% Didascalie compatte e senza rientro sospeso; numeri senza punto finale.
\KOMAoptions{numbers=noendperiod}
\setcapindent{0pt}
\addtokomafont{caption}{\small}
\addtokomafont{captionlabel}{\bfseries}
% Figure e testo sulla stessa pagina quando possibile; pagine di sole figure
% allineate in alto.
\renewcommand{\topfraction}{.85}
\renewcommand{\bottomfraction}{.7}
\renewcommand{\textfraction}{.12}
\renewcommand{\floatpagefraction}{.75}
\makeatletter
\setlength{\@fptop}{0pt}
\makeatother
% Indice: spazio per numeri a due cifre e romani; i titoli lunghi vanno a capo
% prima del numero di pagina (con XeLaTeX un link non si spezza, quindi nell'indice
% è cliccabile solo il numero di pagina; il riquadro del numero accoglie tre cifre).
\PassOptionsToPackage{linktoc=page}{hyperref}
\DeclareTOCStyleEntry[numwidth=2.4em,pagenumberwidth=2.2em,rightindent=3.2em]{tocline}{chapter}
\RedeclareSectionCommand[tocnumwidth=2.6em]{part}
